Every term you are likely to hear, in one line each. Grouped by area rather than
alphabetised, because you will be looking things up by context.

\subsection{Rhythm and timing}\label{th09_glossary-rhythm-and-timing}

\textbf{BPM} Beats per minute. The tempo.

\textbf{Bar} One cycle of the meter. Four beats in almost all dance music.

\textbf{Beat} One quarter note. In house, one kick drum.

\textbf{Step} One sixteenth note. A drum machine bar has sixteen of them.

\textbf{Downbeat} Beat 1 of a bar.

\textbf{Offbeat / upbeat} The "\&" halfway between two beats. Where house hi-hats and basslines
live.

\textbf{Backbeat} Beats 2 and 4. Where the clap or snare goes.

\textbf{Four on the floor} A kick drum on every beat.

\textbf{Two-step} A broken kick pattern, characteristic of UK garage.

\textbf{Half-time} A pattern that implies half the tempo, usually by putting the snare on beat 3
instead of on 2 and 4.

\textbf{Swing / shuffle} Delaying every second subdivision to create a lopsided groove. 50\% is
straight, 66.7\% is a full triplet feel.

\textbf{Quantise} Snapping notes to the nearest grid position.

\textbf{Groove / feel} The overall rhythmic character, including swing, velocity and micro-timing.

\textbf{Syncopation} Accents placed where the meter does not expect them.

\textbf{Polyrhythm} Two rhythmic cycles of different lengths running at once.

\textbf{Ghost note} A deliberately quiet hit between the accented ones.

\textbf{Velocity} How hard a note is played, 0 to 127. Usually maps to volume and brightness.

\textbf{Phrase} A group of bars, usually 4, 8, 16 or 32, forming one unit of the arrangement.

\subsection{Arrangement}\label{th09_glossary-arrangement}

\textbf{Intro} The opening section, usually drums only, designed for the DJ to mix into.

\textbf{Outro} The closing section, same idea in reverse.

\textbf{Breakdown} A section where the drums drop out and the atmosphere takes over.

\textbf{Build / riser section} The tension-increasing section before a drop.

\textbf{Drop} The moment everything returns at full force. The payoff.

\textbf{Riser} A rising noise or tone used to build tension.

\textbf{Impact} A low booming hit marking a drop.

\textbf{Fill} A short variation at the end of a phrase that signals a change is coming.

\textbf{Loop} A repeating section of music.

\textbf{Loop syndrome} The chronic condition of making dozens of loops and finishing no tracks.

\textbf{Stem} A submix of one group of elements (all drums, all vocals) exported separately.

\textbf{Edit / rework / VIP} Versions of a track: a DJ edit for mixing, a rework by someone else,
a VIP as the producer\textquotesingle s own alternate version.

\textbf{Dub / instrumental} A version with the vocal removed or heavily reduced.

\textbf{Extended mix} The full-length club version, as opposed to the radio edit.

\subsection{Harmony}\label{th09_glossary-harmony}

\textbf{Semitone} The smallest step in Western music. A ratio of 2\^{}(1/12).

\textbf{Interval} The distance between two notes, measured in semitones.

\textbf{Scale} The set of notes a piece stays inside.

\textbf{Key} A scale plus a home note.

\textbf{Root} The note a chord is built on.

\textbf{Triad} A three-note chord.

\textbf{Seventh / ninth chord} A chord with a fourth or fifth note stacked on, giving the jazzy,
unresolved quality house is built on.

\textbf{Inversion} The same chord with a different note at the bottom.

\textbf{Voicing} The specific arrangement and spacing of a chord\textquotesingle s notes.

\textbf{Voice leading} Moving between chords by shifting each note as little as possible.

\textbf{Progression} A sequence of chords.

\textbf{Vamp} A short progression, often just two chords, repeated indefinitely.

\textbf{Roman numerals} Key-independent shorthand for chords: i, iv, V and so on.

\textbf{Relative minor} The minor key sharing the same notes as a given major key.

\textbf{Pentatonic} A five-note scale in which nothing sounds wrong.

\textbf{Monophonic / polyphonic} One note at a time versus several. Basses should be monophonic.

\textbf{Camelot wheel} A relabelling of the circle of fifths used by DJs for harmonic mixing.

\subsection{Synthesis}\label{th09_glossary-synthesis}

\textbf{Oscillator} The part that generates the raw tone.

\textbf{Waveform} The oscillator\textquotesingle s shape: sine, triangle, square, sawtooth, noise.

\textbf{Subtractive synthesis} Starting with a rich waveform and filtering harmonics away. The
dominant method.

\textbf{Additive synthesis} Building a sound by summing individual sine partials.

\textbf{FM synthesis} Using one oscillator to modulate another\textquotesingle s frequency at audio rate.

\textbf{Wavetable synthesis} Scanning through stored waveforms so the timbre morphs over time.

\textbf{Filter} Removes frequencies above or below a cutoff.

\textbf{Cutoff} The filter\textquotesingle s corner frequency.

\textbf{Resonance} A boost right at the cutoff. High resonance is the acid sound.

\textbf{Slope} How steeply a filter cuts, in dB per octave. 12, 18 and 24 are common.

\textbf{Envelope / ADSR} Attack, Decay, Sustain, Release. How a parameter changes over one note.
Sustain is a level; the other three are times.

\textbf{LFO} A low frequency oscillator used as a controller rather than a sound source.

\textbf{Modulation} Any parameter controlling another parameter.

\textbf{Portamento / glide / slide} Sliding continuously from one pitch to the next.

\textbf{Unison / detune} Stacking several slightly out-of-tune copies of an oscillator for
thickness. The supersaw.

\textbf{Patch / preset} A saved synth setting.

\textbf{Polyphony} How many notes a synth can play at once.

\subsection{Sampling}\label{th09_glossary-sampling}

\textbf{One-shot} A single sound, triggered whole.

\textbf{Sample pack} A commercial collection of one-shots and loops.

\textbf{Chop} Cutting a sample into pieces and resequencing them.

\textbf{Slicing} Automatically dividing a loop at its transients.

\textbf{Time stretch} Changing length without changing pitch.

\textbf{Pitch shift} Changing pitch without changing length.

\textbf{Vocal chop} A short vocal fragment played as an instrument.

\textbf{Break} A drum-only section of an old record, sampled. The Amen break is the famous one.

\textbf{Clearance} Legal permission to use a sample from a copyrighted recording.

\textbf{Royalty free} Sample material licensed for use in your own productions.

\subsection{Mixing and mastering}\label{th09_glossary-mixing-and-mastering}

\textbf{Gain staging} Setting sensible levels through the whole signal chain.

\textbf{dBFS} Decibels relative to full scale. 0 dBFS is the digital ceiling.

\textbf{Headroom} How much space is left before clipping.

\textbf{Clipping} Exceeding the ceiling, producing distortion.

\textbf{EQ} Equalisation. Boosting or cutting specific frequency ranges.

\textbf{Highpass / lowpass} Filters removing lows or highs respectively.

\textbf{Compression} Reducing dynamic range so the whole thing can be louder.

\textbf{Threshold, ratio, attack, release} The four main compressor controls.

\textbf{Sidechain} Using one signal to control a processor on another. Usually the kick ducking
the bass.

\textbf{Limiter} A compressor with a very high ratio, used to catch peaks.

\textbf{Send / return / bus} Routing a copy of several channels to a shared effect.

\textbf{Dry / wet} Unprocessed versus processed signal.

\textbf{Pan} Position in the stereo field.

\textbf{Mono compatibility} Whether the mix survives being summed to a single channel.

\textbf{Phase} The relative timing of waveforms. Misaligned phase causes cancellation.

\textbf{LUFS} The standard perceptual loudness measurement.

\textbf{True peak (dBTP)} Peak level accounting for what happens between samples. Keep at or
below −1.

\textbf{Loudness war} The long-running trend of mastering everything as loud as possible.

\textbf{Mastering} The final polishing and loudness stage.

\textbf{Reference track} A commercial release you A/B against while mixing.

\subsection{Software and hardware}\label{th09_glossary-software-and-hardware}

\textbf{DAW} Digital audio workstation. FL Studio, Ableton Live, Logic, Bitwig.

\textbf{Plugin} Software instrument or effect loaded inside a DAW.

\textbf{VST / AU / CLAP} Plugin formats. VST is the Windows standard.

\textbf{MIDI} A protocol for note and control data. Not audio.

\textbf{Piano roll} The grid editor where MIDI notes are drawn.

\textbf{Step sequencer} A grid of on/off buttons for programming drums.

\textbf{Automation} Recorded parameter changes over time.

\textbf{ASIO} The low-latency Windows audio driver standard. Required for usable performance.

\textbf{Buffer size} How much audio is processed at once. Smaller means lower latency and higher
CPU load.

\textbf{Latency} The delay between an action and hearing the result.

\textbf{Bounce / render / export} Writing the project out as an audio file.

\textbf{Freeze} Rendering a track to audio temporarily to save CPU.

\textbf{Controller} A hardware device sending MIDI to your DAW.

\textbf{Audio interface} External hardware for converting audio in and out of the computer.

\subsection{Machines you will hear named}\label{th09_glossary-machines-you-will-hear-named}

\textbf{TR-909} Roland drum machine, 1983. The house and techno standard.

\textbf{TR-808} Roland drum machine, 1980. Deep booming kick, cowbell. Hip-hop and electro.

\textbf{TR-707} Roland drum machine, 1985. Cheaper, drier, all over early Chicago records.

\textbf{TB-303} Roland bass synth, 1981. A commercial failure that invented acid house.

\textbf{Juno-60 / Juno-106} Roland polysynths with a famous chorus. The deep house pad sound.

\textbf{SP-1200 / MPC} Samplers. The foundation of hip-hop and a lot of house.

\textbf{Technics SL-1200} The DJ turntable.

\textbf{CDJ} Pioneer\textquotesingle s standard club media player.

\subsection{Scene and culture}\label{th09_glossary-scene-and-culture}

\textbf{House} The genre, named after the Warehouse in Chicago.

\textbf{Jack} To dance to house. Also a jack track: a stripped, hypnotic, drum-led record.

\textbf{Resident} A DJ with a regular slot at one club.

\textbf{Set} A DJ\textquotesingle s performance.

\textbf{Beatmatching} Aligning the tempo and phase of two records.

\textbf{Mixing in key} Selecting records with compatible keys.

\textbf{Warm-up} The early part of the night, played slower and deeper. A specialised skill.

\textbf{Peak time} The busiest, hardest part of the night.

\textbf{White label / promo} A limited pressing without full artwork, usually for DJs.

\textbf{Crate digging} Searching for records in shops.

\textbf{Rave} An event, historically illegal or unlicensed.

\textbf{Bootleg / mashup} An unofficial edit combining or reworking existing records.

\textbf{Four to the floor} See four on the floor. Same thing, said differently, by everyone.
